\subsection{Сценарий 2: near-real-time-витрина при потоковой записи}

В сценарии сравнивались форматы ORC, Parquet и Vortex на таблице с первичным ключом
(схема журнала событий из 25 столбцов, таблица~\ref{tab:schema}) при потоковой записи
с захватом изменений и параллельных аналитических запросах. Формат Lance в
предварительном прогоне этого сценария оказался стабильно худшим по времени чтения
(в \num{2}--\num{4} раза медленнее ORC и Parquet) и из итогового сравнения исключён
(см. подраздел~\ref{subsec:file-formats}). Результаты получены в
представительном запуске (три повтора полного цикла, ограничение скорости фазы
обновления --- порядка 100\,000~строк/с, объём фазы обновления --- около
\num{200000000} строк, интервал чекпойнтов --- 240~с).

Пропускная способность фазы начальной загрузки (последовательные
ключи, реальная запись на диск) практически одинакова у всех трёх форматов
(таблица~\ref{tab:write}); различие в пределах разброса. В фазе обновления показатели
также близки, но интерпретируются с оговоркой: при постоянном значении поля
последовательности значительная доля обновлений дедуплицируется в буфере записи и не
доходит до хранилища, поэтому величина отражает скорость генерации и дедупликации в
памяти. Вывод: Vortex не уступает Parquet и ORC по пропускной способности записи ---
вторая гипотеза работы подтверждается.

\begin{table}
    \caption{Пропускная способность записи, тыс. строк в секунду (медиана по повторам)}
    \label{tab:write}
    \begin{tabularx}{\textwidth}{|X|c|c|}\hline
    \textbf{Формат} & \textbf{Начальная загрузка (INIT)} & \textbf{Обновление (UPDATE)$^{*}$} \\ \hline
    ORC     & 1505 & 89 \\ \hline
    Parquet & 1504 & 99 \\ \hline
    Vortex  & 1499 & 89 \\ \hline
    \multicolumn{3}{|p{0.95\textwidth}|}{\footnotesize $^{*}$~значение фазы обновления отражает скорость генерации и дедупликации в буфере записи (постоянное поле последовательности), а не скорость записи на диск.} \\ \hline
    \end{tabularx}
\end{table}

Помимо времени запросов, форматы различаются по занимаемому объёму хранилища.
Размер таблицы после полного цикла записи приведён в таблице~\ref{tab:nrt-size}. ORC даёт самый компактный результат, Parquet --- на
\SI{6}{\percent} больше, Vortex --- на \SI{18}{\percent} больше ORC. Различие в
пределах нескольких процентов сопоставимо с разбросом и не влияет на стоимость
хранения принципиально, однако заслуживает комментария: на TPC-DS Vortex и Parquet
были эквивалентны по объёму, а на потоковой схеме Vortex несколько менее эффективен.
Возможные причины: высокая доля null-значений в столбцах схемы
(\texttt{parent\_req\_id} --- \SI{30}{\percent}, \texttt{ad\_id} --- \SI{50}{\percent},
\texttt{revenue} --- \SI{90}{\percent}) и случайные идентификаторы в UPDATE-фазе
снижают плотность FSST-кодирования строк-префиксов; кроме того, фрагментация LSM
на множество мелких файлов уровня~0 увеличивает долю накладных расходов (футер,
схема) от общего объёма (у Vortex заметно больше объектов в хранилище: 1883 против
1731--1733 у Parquet/ORC).

\begin{table}
    \caption{Объём хранилища таблицы near-real-time-витрины после полного цикла записи}
    \label{tab:nrt-size}
    \begin{tabularx}{\textwidth}{|X|c|c|c|}\hline
    \textbf{Формат} & \textbf{Объём, ГиБ} & \textbf{Число объектов} & \textbf{Отношение к ORC} \\ \hline
    ORC     & 17 & 1731 & ${\times}1.00$ \\ \hline
    Parquet & 18 & 1733 & ${\times}1.06$ \\ \hline
    Vortex  & 20 & 1883 & ${\times}1.18$ \\ \hline
    \end{tabularx}
\end{table}

Выбор основной метрики усреднения требует пояснения. В DURING-фазе отдельные сэмплы чтения
могут попадать на момент массового сброса данных по чекпойнту, когда writer
одновременно пишет в объектное хранилище и данные таблицы, и состояние чекпойнта
Flink (последнее размещается в той же директории объектного хранилища). Такое
совпадение даёт единичные «хвостовые» сэмплы длительностью в десятки секунд;
среднее по сэмплам, особенно при небольшом числе повторов, чувствительно к таким
попаданиям. Поэтому в качестве основной метрики далее используется медиана
по выборкам и повторам --- она отражает типичную задержку чтения, к которой и
чувствителен потребитель витрины данных, и устойчива к единичным выбросам.
Среднее приводится отдельно как индикатор подверженности формата выбросам.

Сводные результаты по типичной задержке чтения (медиане) приведены в таблице~\ref{tab:read-med}. По медиане Vortex --- быстрейший формат в обеих фазах:
в фазе FINAL отрыв заметный ($\approx 23\,\%$ относительно ORC и Parquet), в фазе
DURING --- скромный (доли процента). При этом ORC и Parquet по медиане практически
не меняют времени между фазами DURING и FINAL, а Vortex в фазе FINAL заметно
быстрее самого себя в фазе DURING ($\approx 22\,\%$ улучшения) --- эффект уплотнения
данных: после остановки записи и работы фонового слияния файлы сортированы по
ключу, и специализированные кодеки Vortex (прежде всего FSST на строках-префиксах)
работают в полной плотности.

\begin{table}
    \caption{Время выполнения запроса, мс (по всем запросам, выборкам и повторам): медиана и среднее}
    \label{tab:read-med}
    \begin{tabularx}{\textwidth}{|X|c|c|c|c|}\hline
    \multirow{2}{*}{\textbf{Формат}} & \multicolumn{2}{c|}{\textbf{Фаза DURING}} & \multicolumn{2}{c|}{\textbf{Фаза FINAL}} \\ \hhline{~----}
    & медиана & среднее & медиана & среднее \\ \hline
    ORC     & 7756 & 7594 & 7767 & 7633 \\ \hline
    Parquet & 7824 & 9481 & 7797 & 7767 \\ \hline
    Vortex  & 7677 & 7712 & 6004 & 6736 \\ \hline
    \end{tabularx}
\end{table}

Подробные результаты по медиане в фазе FINAL приведены в таблице~\ref{tab:read-final}. Vortex
выигрывает на девяти запросах из тринадцати; ORC --- на двух
(\texttt{AGG\_COUNT\_DISTINCT}, \texttt{AGG\_SUM\_BYTES}, разница в пределах
единиц процентов); Parquet --- на двух (\texttt{AGG\_GROUP\_COUNTRY},
\texttt{AGG\_MULTIDIM}). Наибольший отрыв Vortex --- на запросах, чувствительных к
упорядоченности данных и к качеству кодирования строк с регулярной структурой:
\texttt{FUNNEL} (многошаговый запрос) --- $\approx 33\,\%$,
\texttt{JOIN\_REVENUE\_PROFILE} --- $\approx 20\,\%$, \texttt{AGG\_FILTERED}
(подсчёт с предикатом) --- $\approx 21\,\%$, \texttt{RANGE\_SESSION} (диапазонный
фильтр по идентификатору сессии) --- $\approx 0.2\,\%$ (медианы почти равны, но
среднее у Vortex заметно ниже).

\begin{table}
    \caption{Медиана времени выполнения запросов в фазе FINAL, мс}
    \label{tab:read-final}
    \begin{tabularx}{\textwidth}{|X|c|c|c|}\hline
    \textbf{Запрос} & \textbf{ORC} & \textbf{Parquet} & \textbf{Vortex} \\ \hline
    \texttt{POINT}                  & 3841  & 4169  & 3780 \\ \hline
    \texttt{FULLSCAN\_PROJ}          & 5888  & 5839  & 5829 \\ \hline
    \texttt{RANGE\_SESSION}          & 5837  & 7851  & 5826 \\ \hline
    \texttt{RANGE\_URL}              & 5823  & 5873  & 5787 \\ \hline
    \texttt{AGG\_SUM\_BYTES}         & 5703  & 5868  & 5808 \\ \hline
    \texttt{AGG\_COUNT\_DISTINCT}    & 7823  & 9870  & 7842 \\ \hline
    \texttt{AGG\_FILTERED}           & 9805  & 9874  & 7783 \\ \hline
    \texttt{AGG\_GROUP\_COUNTRY}     & 5804  & 5778  & 5825 \\ \hline
    \texttt{AGG\_MULTIDIM}           & 7863  & 6407  & 7922 \\ \hline
    \texttt{FUNNEL}                  & 12303 & 12162 & 8221 \\ \hline
    \texttt{JOIN\_USER\_ACTIVITY}    & 8134  & 10120 & 8044 \\ \hline
    \texttt{JOIN\_SESSION\_LATENCY}  & 10077 & 8103  & 7992 \\ \hline
    \texttt{JOIN\_REVENUE\_PROFILE}  & 7709  & 8147  & 6195 \\ \hline
    \multicolumn{1}{|r|}{выигрышей} & 2 & 2 & 9 \\ \hline
    \end{tabularx}
\end{table}

Подробные результаты в фазе DURING приведены в таблице~\ref{tab:read-during}. По медиане Vortex также выигрывает большинство
запросов (девять из тринадцати). Наибольший отрыв --- на запросах, где работает
проталкивание предиката и плотное кодирование строк с регулярной структурой:
\texttt{RANGE\_SESSION} --- $\approx 25\,\%$ быстрее ORC и в три раза быстрее
Parquet (Parquet на этом запросе под параллельной записью ловит крупный выброс ---
медиана 17\,684~мс), \texttt{FUNNEL} --- $\approx 15\,\%$, \texttt{RANGE\_URL} ---
$\approx 4\,\%$. ORC выигрывает на простой группировке и многомерной агрегации
(\texttt{AGG\_GROUP\_COUNTRY}, \texttt{AGG\_MULTIDIM}, разница в пределах разброса);
Parquet --- на точечном запросе и полном сканировании с проекцией
(\texttt{POINT}, \texttt{FULLSCAN\_PROJ}, разница в пределах долей процента).

\begin{table}
    \caption{Медиана времени выполнения запросов в фазе DURING, мс}
    \label{tab:read-during}
    \begin{tabularx}{\textwidth}{|X|c|c|c|}\hline
    \textbf{Запрос} & \textbf{ORC} & \textbf{Parquet} & \textbf{Vortex} \\ \hline
    \texttt{POINT}                  & 3724  & 3706  & 3721 \\ \hline
    \texttt{FULLSCAN\_PROJ}          & 5845  & 5773  & 5812 \\ \hline
    \texttt{RANGE\_SESSION}          & 7800  & 17684 & 5863 \\ \hline
    \texttt{RANGE\_URL}              & 5932  & 6133  & 5718 \\ \hline
    \texttt{AGG\_SUM\_BYTES}         & 5716  & 5746  & 5672 \\ \hline
    \texttt{AGG\_COUNT\_DISTINCT}    & 7877  & 13700 & 7755 \\ \hline
    \texttt{AGG\_FILTERED}           & 8035  & 9789  & 7679 \\ \hline
    \texttt{AGG\_GROUP\_COUNTRY}     & 5776  & 5776  & 5803 \\ \hline
    \texttt{AGG\_MULTIDIM}           & 5926  & 7696  & 5953 \\ \hline
    \texttt{FUNNEL}                  & 12005 & 11999 & 10193 \\ \hline
    \texttt{JOIN\_USER\_ACTIVITY}    & 7998  & 8094  & 7971 \\ \hline
    \texttt{JOIN\_SESSION\_LATENCY}  & 8024  & 8091  & 8006 \\ \hline
    \texttt{JOIN\_REVENUE\_PROFILE}  & 8097  & 8179  & 8085 \\ \hline
    \multicolumn{1}{|r|}{выигрышей} & 2 & 2 & 9 \\ \hline
    \end{tabularx}
\end{table}

Соотнесение среднего и
медианы по DURING-фазе полезно для оценки тяжести хвостов задержки. Для ORC
среднее (7594~мс) и медиана (7756~мс) практически совпадают, для Vortex среднее
(7712~мс) и медиана (7677~мс) также совпадают --- это означает, что у обоих
форматов крупных выбросов задержки в этом запуске нет. Для Parquet среднее
(9481~мс) на $\approx 21\,\%$ выше медианы (7824~мс) --- среднее инфлировано
несколькими сэмплами в десятки секунд (например, на запросе
\texttt{AGG\_COUNT\_DISTINCT} одна из выборок заняла 24\,801~мс, тогда как
медиана по этому запросу --- 13\,700~мс). Это симптом, согласующийся с
архитектурным предположением о пуле соединений: читатель Parquet ходит в
объектное хранилище через общий пул слоя совместимости Hadoop S3A и в момент
массового сброса данных по чекпойнту иногда простаивает в ожидании соединения,
давая катастрофически медленную выборку; читатель Vortex использует независимый
нативный пул и таких выбросов не даёт. Читатель ORC, формально использующий тот
же общий пул, в этом запуске также не даёт выбросов --- благодаря структурным
особенностям формата: крупные полосы (stripes, типично около \SI{64}{\mega\byte})
и эффективное проталкивание предикатов на уровне полос и индексных групп по
\num{10000} строк сокращают число обращений к хранилищу за один скан, поэтому
читатель ORC реже простаивает в очереди за соединением, даже когда общий пул под
нагрузкой.

Промежуточные выводы по сценарию таковы. Vortex эквивалентен Parquet и ORC по
пропускной способности записи (гипотеза~2 подтверждена). По типичной задержке
чтения (медиана) Vortex --- быстрейший формат в обеих фазах и на большинстве
запросов (девять из тринадцати в каждой фазе); особенно выраженно --- на запросах
с диапазонным фильтром по строке-префиксу, селективным предикатом перед агрегацией
и в многошаговой аналитике. На простой группировке и многомерной агрегации без
предиката Vortex незначительно (в пределах единиц процентов) уступает конкурентам
--- следствие отсутствия проталкивания агрегатов в интерфейсе форматов Apache
Paimon (гипотеза~3 подтверждается). Между фазами FINAL и DURING медианы ORC и
Parquet практически не меняются, медиана Vortex в DURING выше своего же FINAL на
$\approx 28\,\%$: специализированные кодеки Vortex (прежде всего FSST) дают
наибольшую плотность на упорядоченных данных, а в свежих файлах уровня~0,
сбрасываемых по чекпойнту, порядок ещё не уплотнён слиянием. По хвостам задержки
ORC и Vortex показывают близкие, узкие распределения (среднее почти равно медиане),
Parquet --- более широкое распределение с эпизодическими выбросами в десятки
секунд в моменты сброса данных по чекпойнту. Это согласуется с архитектурным
предположением о выигрыше Vortex за счёт независимого нативного пула соединений с
объектным хранилищем по сравнению с Parquet (гипотеза~4 подтверждается для пары
Vortex--Parquet); ORC достигает аналогично узких хвостов иначе --- меньшим
числом обращений к хранилищу за один скан за счёт крупных полос (stripes) и
гранулярного пропуска данных по индексным группам.
